\section{Discussion and Open Questions}
\label{sec:discussion}

\subsection{What this paper establishes}

\paragraph{Direct $\volR$ measurement is privacy-incompatible.}
The framework's commitments
in~\citep{lasser2026exo,lasser2026phase,lasser2026wamura} rule
out the panopticon-style observation that direct measurement
would require.  This is a structural finding, not a practical
limitation.

\paragraph{Revealed sacrifice is a privacy-minimal surrogate.}
The sacrifice channel satisfies three formal privacy properties
(Theorem~\ref{thm:privacy}) plus two institutional preconditions
(Proposition~\ref{prop:privacy_institutional}), and composes
with the commitment layer of~\citep{lasser2026exo}
(Propositions~\ref{prop:commit_preserves}
and~\ref{prop:zk_aggregation}) to produce a privacy-minimal
observation mechanism on distinct-category batches.

\paragraph{The B-to-C gap is characterized by trade data.}
The aggregate lower bound
(Theorem~\ref{thm:aggregate_bound}) lower-bounds the exercised
fraction $\beta$ (equivalently, upper-bounds the B-to-C gap
$1 - \beta$) on the observed subset, under S0--S2 and S4(a).
The constructed lower bound is non-decreasing in observed
admissible events: Part~A via sub-batch supremum
(Proposition~\ref{prop:monotone_part_a}) and Part~B via
per-capability max-attribution
(Proposition~\ref{prop:monotone}).

\paragraph{Two sacrifice channels, substitutable with residuals.}
The money and time channels are substitutable via cooperative
markets (employment and hiring) and triangulate the same
underlying sacrifice substrate.  The joint aggregate weakly
dominates either channel alone unconditionally; strict
domination obtains at the non-fungibility residuals
(agent-specific time, capital-concentrated access) when the
trade window contains an event in the relevant residual
category, per Corollary~\ref{cor:joint_coverage}.

\paragraph{$\volR$ is not self-balancing.}
$\volR$ exhibits non-invariance under capability removal:
eliminating an unexercised capability leaves $\volR$ unchanged
while $\volL$ contracts
(Theorem~\ref{thm:axiom_inheritance}, M6(a); Example
\ref{ex:volR_removal_invariance}).  The self-balancing
protection $\volL$ enjoys---every removal contracts the
measure---therefore does not transfer to $\volR$
(Corollary~\ref{cor:not_self_balancing}).  $\volR$ is a
diagnostic, not an objective.  This structural finding is
independent of the observation mechanism.

\paragraph{Complementarity with controlled relaxation.}
Together, this paper and~\citep{lasser2026wamura} cover the
risk and value dimensions of the
preemptive-restriction criterion of~\citep[Proposition~1]{lasser2026horizon}.
Neither is sufficient alone.

\subsection{Limitations}

\paragraph{Assumption load.}
The full per-capability refinement
(Theorem~\ref{thm:aggregate_bound} Part~B) requires S0--S5 plus
within-bundle poset-independence
(Definition~\ref{def:bundle_decomp}).  S0 (calibration) and S1
(free choice) are the most empirically demanding, and neither
is ZK-verifiable (Remark~\ref{rem:zk_limits}).  Part~A (S0--S2
and S4(a) only) is the more defensible result; Part~B should be
understood as the refinement available under stronger
conditions.

\paragraph{B-to-C gap inheritance.}
The damage bound of~\citep[Theorem~1]{lasser2026wamura} bounds
contraction of $\volL$, not of the experiential target
directly.  The present paper provides a lower bound on $\volR$
as a diagnostic, but the diagnostic is itself subject to the
B-to-C gap on the benchmarked side: the sacrifice signal
$\Delta\volL(X)$ is only as accurate as $\volL$'s measurement
of the benchmarked capability.

\paragraph{External admissibility layer.}
Several non-ZK-verifiable hypotheses require institutional
attestation to hold operationally: S0 (valuation bounded by
exercise), S1 (free choice and non-degradation), S4(a)
(valuation additivity), and the genuine-exchange conditions
(iii)--(iv) of
Definition~\ref{def:sacrifice_event}.
Remark~\ref{rem:external_attestation} describes the operational
shape of this attestation layer.  The specific filter for S1's
non-degradation condition---rejecting below-sufficiency
sacrifices as structurally involuntary---is developed in the
companion paper~\citep{lasser2026paper8b} and treated here as
an external precondition.

\subsection{Open questions}

\begin{enumerate}

\item \label{oq:hedonic}
\textbf{Hedonic disaggregation of bundles.}
Theorem~\ref{thm:aggregate_bound} Part~B requires event-local
decomposition coefficients $\alpha_{n,c}$.  Classical hedonic
regression~\citep{rosen1974hedonic} handles the money channel.
The time-sacrifice channel needs an analogous framework: how
does one decompose a block of time spent on a complex activity
(parenting, which involves teaching, emotional regulation,
physical care, relationship-building) into per-capability
contributions?

\item \label{oq:nonstationarity}
\textbf{Non-stationarity of revealed preferences.}
Preferences shift over time.  The lower bound from past trades
may overstate $\volR$ at present.  An EWMA decay paralleling
the risk-trust dynamics of~\citep[Definition~4]{lasser2026scm}
seems natural but requires empirical calibration.

\item \label{oq:coercion}
\textbf{Coerced sacrifice.}
Assumption S1 (free choice) is load-bearing.  Under duress,
the agent's trade does not reveal preferences.  Filtering
coerced sacrifice events may require commitment-based
attestation of choice sets~\citep{lasser2026exo}.

\item \label{oq:market_failure}
\textbf{Market-failure cases.}
In contexts without functioning markets (pre-capitalist
societies, closed economies, intra-family sharing), the
money-sacrifice channel is quiet even though $\volR$ activity
is real.  Does the time-sacrifice channel alone suffice, or
does the framework need a non-market sacrifice variant?

\item \label{oq:commitment_infra}
\textbf{Commitment infrastructure at scale.}
A fully committed trade ledger at the scale of a modern
economy is an infrastructure problem.  The privacy guarantee
(Section~\ref{sec:commitment}) is conditional on the
commitment infrastructure existing, which is the remit
of~\citep{lasser2026exo}.

\item \label{oq:calibration}
\textbf{Population-level calibration (S0).}
Assumption S0 requires that individual agents do not overvalue
unbenchmarked bundles.  A robust version of
Theorem~\ref{thm:aggregate_bound} would replace S0 with a
weaker population-level condition:
$E[\Ui(Y)] \leq \volR(Y) + \epsilon$ for some bounded bias
$\epsilon$.  Specific behavioural mechanisms that stress
individual-agent S0 and motivate a population-level weakening
include: (i) \emph{ritualized payment}, where a transaction's
nominal price is symbolic rather than exchange-valued (a token
fee to ``soften'' receiving something for free);
(ii) \emph{gift aversion}, where agents refuse free offers
because they interpret absence of price as signalling hidden
defect, inflating observed sacrifices on nominally free goods;
(iii) \emph{time-for-money inefficiency}, where a given hour's
wage rate $r_i$ does not reflect the agent's marginal valuation
(labour-market frictions, status wages, family-care default
allocations).  Each of these produces systematic biases on
observed sacrifice magnitudes that the per-agent S0 cannot
absorb; a population-level reformulation that lets such biases
cancel (or bounds the residual) is the natural follow-up.

\end{enumerate}

\paragraph{Further open questions in the companion
paper.} Open questions on threshold governance, need
identification, capability-stuffing as structural avoidance,
and the full architecture of gap-decomposition governance are
collected in the companion paper~\citep{lasser2026paper8b}.
